\subsection{A degree-two defect at every prime in a deep resolved tail}

\begin{theorem}[A-P01: Proved; deep coefficient resolutions can create low-degree torsion at every prime]
\label{AP01}
Let $p$ be any prime, $R=\mathbb Z_{(p)}$, $S=\Gamma_R(T)$,
and $I_1=\bigoplus_{w\geq1}S^w$.
For every $b\geq3$ put $s=pb$ and start with the actual
strict simplicial PD algebra
\[
 X=G_R\otimes_R\operatorname{FreePD}_R
          (x_0K(2)\xrightarrow{d}y_0K(2)),
 \qquad |x_0|=1,\quad |y_0|=2,\quad
 \operatorname{wt}(x_0)=\operatorname{wt}(y_0)=b.
 \tag{DP.1}
\]
The element $z_0=\gamma_p(y_0)$ in simplex two has exact
order $p$ in $H^{2p}(X_2)$, and
$\eta_0=x_0\gamma_{p-1}(y_0)$ satisfies $d\eta_0=pz_0$.
We construct below a free coefficient resolution of
$I_1\langle s\rangle$ and its nonnegative resolved-tail
PD attachment $Y^{\rm tf}$ for this pair.
It has an actual projection to the admissible theory $G_R$,
coefficient $H^0$, rational coefficient comparison,
and unchanged simplices zero and one.

After the entire iterated based PD saturation of A-SAT01
at cutoff $s+1$, its simplex-two cohomology contains a
class of exact order $p$ in form degree two and weight
\[
                        2s+2p(p-1)+1.               \tag{DP.2}
\]
The statement holds for arbitrarily large $p$, although
the class being repaired originally has degree $2p$.
It is a defect of this actual resolved-tail construction;
it is not a counterexample to common admissible sources.
The proof also gives an actual next degree-two repair of
the displayed class with the chosen projection to $G_R$.
\end{theorem}

\begin{proof}
We specify the coefficient resolution completely in the
range used, and give an explicit rule for its extension
outside that range. Set
\[
 t=T,\qquad u=T^{[p]},\qquad c=(p-1)!,\qquad t^p=pcu.
 \tag{DP.3}
\]
The ring $B=R[t,u]/(t^p-pcu)$ injects into $S$.
Indeed it is free over $R[u]$ on $1,t,\ldots,t^{p-1}$,
and these monomials have distinct weights and nonzero
images in $\mathbb Q[T]$.
In weights below $p^2$ the inclusion is an isomorphism:
for $0\leq r<p$ and $0\leq j<p$, the coefficient of
$T^{[jp+r]}$ in $t^ru^j$ is a $p$-adic unit by Legendre's
valuation formula. At weight $p^2$ the new generator
$v=T^{[p^2]}$ satisfies
\[
 u^p=pc'v,\qquad
 c'=\frac{(p^2)!}{p(p!)^p}\in R^\times .            \tag{DP.4}
\]
The two coefficient valuations in question are respectively
zero and one; the formulas also cover $p=2$.

Begin with free generators $e_1,e_p$ augmenting to $t,u$.
The following two matrices alternate as differentials
of a free resolution of $(t,u)$ over $B$:
\[
 M=\begin{pmatrix}t^{p-1}&u\\-pc&-t\end{pmatrix},
 \qquad
 N=\begin{pmatrix}t&u\\-pc&-t^{p-1}\end{pmatrix}.
 \tag{DP.5}
\]
The first is $Q_1\to Q_0$, the second is $Q_2\to Q_1$.
We prove their exactness directly. Modulo $p$, the displayed
monomial basis identifies $B/p$ with
$\mathbb F_p[t,u]/(t^p)$. Consequently
\[
 (pc:t)=(pc,t^{p-1}),\qquad
 (pc:t^{p-1})=(pc,t).                              \tag{DP.6}
\]
The kernel of $Q_0\to(t,u)$ is generated by
$(t^{p-1},-pc)$ and $(u,-t)$: an equation
$ta+ub=0$, multiplied by $t^{p-1}$ and then cancelled
by the nonzero element $u$ in the domain $B$, gives
$pc\,a=-t^{p-1}b$, so the second
colon identity expresses $b=-pc\,h-tk$ and then
$a=t^{p-1}h+uk$. The injectivity of multiplication by
$pc$ justifies this last equality.
For an element in $\ker M$, its second coordinate
requires $pc\,h=-tk$. The first colon identity writes
$k=-pc\,a-t^{p-1}b$ and gives $h=ta+ub$; these are
exactly the columns of $N$. The other equation follows
from (DP.3). The same argument with the second colon
identity proves $\ker N=\operatorname{im}M$.
This proves all exactness assertions, including every
stage of the alternating resolution.

Its two generator weights are
\[
 Q_{2j}:\ (jp+1,(j+1)p),\qquad
 Q_{2j+1}:\ ((j+1)p,(j+1)p+1).                     \tag{DP.7}
\]
To resolve the full ideal $I_1$ over $S$, add
$e_{p^2}$ in degree zero and weight $p^2$, and in degree
one add
\[
                 r_{p^2}=u^{p-1}e_p-pc'e_{p^2}.
 \tag{DP.8}
\]
These additions give every new kernel element at that
weight. In fact its $e_{p^2}$ coordinate must be divisible
by $p$, since the old coefficient image in weight $p^2$
is precisely $pR v$; subtracting a multiple of (DP.8)
reduces to the old kernel. In a relation among degree-one
generators of weight $p^2$, the $e_{p^2}$ coordinate is
$-pc'$ times the coefficient of (DP.8). Torsionfreeness
forces that coefficient to be zero. Thus no new degree-two
syzygy appears at or below weight $p^2$.
All higher kernels in that range are consequently exactly
those of (DP.5). The coefficient ring, in every entry that
occurs, is the already identified low-weight part of $B$.

For clarity, extend this to an actual full resolution as
follows. In degree zero also include $e_{p^i}$ for every
$i\geq3$; the usual base-$p$ monomial expression, with
unit coefficient, generates every positive-weight element
of $S$ from these generators and $e_1,e_p,e_{p^2}$.
At each resolution stage adjoin a free generator for every
homogeneous kernel element of weight greater than $p^2$,
with that element as its differential. Retain the listed
low-weight generators and differentials, which already
surject onto the whole kernel in that range.
Nonnegative coefficient weights prevent a new high-weight
generator from changing any of the verified low-weight
groups. This recursive rule gives an actual weight-graded
free resolution $Q_\bullet\twoheadrightarrow I_1$ with
exactly (DP.5)--(DP.8) in the required range.

Put $q=2p$. On shifted copies of this resolution form
\[
 P^j=Q_{q-1-j}^{u}\langle s\rangle
               \oplus Q_{q-2-j}^{v}\langle s\rangle,
 \quad
 d(u_i)=-u_{\partial i},\qquad
 d(v_i)=p u_i+v_{\partial i},                       \tag{DP.9}
\]
with negative-index $Q$ equal to zero.
Its only cohomology is $I_1\langle s\rangle/p$ in
degree $q-1$. Define the actual nonnegative complex by
\[
 W^0=P^0/dP^{-1},\qquad W^j=P^j\quad(j>0).
 \tag{DP.10}
\]
The differential injects $W^0$ into $P^1$ and identifies
it with $\ker(P^1\to P^2)$, because $H^0(P)=H^1(P)=0$.
In particular $W$ is torsionfree and $H^0(W)=H^1(W)=0$.
Over $\mathbb Q$, the homotopy $u_i\mapsto v_i/p$,
$v_i\mapsto0$ contracts (DP.9). Its quotient on degree
zero and its unchanged positive components contract
(DP.10). This specifies the rational comparison actually
needed below.

Give every $W$ generator arity one and set
$E=\operatorname{im}(\operatorname{FreePD}_S(W)
 \to\operatorname{FreePD}_{S_{\mathbb Q}}(W_{\mathbb Q}))$.
Use the following actual generators, with the stated
unshifted weights; every displayed new generator also
has the common weight shift $s$:
\[
\begin{array}{c|c|c}
\text{generator}&\text{resolution position}&\text{unshifted weight}\\\hline
 C&v(Q_{2p-2},\text{first})&p(p-1)+1\\
 K&v(Q_{2p-2},\text{second})&p^2\\
 H&u(Q_{2p-1},\text{first})&p^2\\
 x&v(Q_{2p-3},\text{first})&p(p-1)\\
 y&v(Q_{2p-3},\text{second})&p(p-1)+1\\
 g&u(Q_{2p-2},\text{first})&p(p-1)+1\\
 r&u(Q_{2p-3},\text{first})&p(p-1)\\
 a&v(Q_{2p-4},\text{first})&p(p-2)+1\\
 b_1&v(Q_{2p-4},\text{second})&p(p-1)
\end{array}                                                       \tag{DP.11}
\]
The first three are zero-forms, $x,y,g$ are one-forms,
and $r,a,b_1$ are two-forms. The quotient relation from
$v(Q_{2p-1},\text{first})$ and the adjacent differentials
are exactly
\[
 \begin{gathered}
 t^{p-1}C=p(cK-H),\\
 dC=pg+tx-pcy,\qquad dx=pr+t^{p-1}a-pc b_1.
 \end{gathered}                                                    \tag{DP.12}
\]
These identities use the complete coefficient resolution,
including (DP.8), rather than a deletion of its generator.

Define the actual integral two-form
\[
             \xi=gx-cyx+Cr-cCb_1+(cK-H)a.
 \tag{DP.13}
\]
The alternating identity $x^2=0$ and (DP.12) give
\[
 \begin{aligned}
 d(Cx)&=(pg+tx-pcy)x+C(pr+t^{p-1}a-pcb_1)\\
      &=p\xi.
 \end{aligned}                                                    \tag{DP.14}
\]
Thus $d\xi=0$ in the torsionfree algebra.
Its rational exterior monomial $gx$ has coefficient one
and uses two independent free positive-form generators.
The other terms have different monomials or form-degree
factors, so $\xi$ is nonzero.

Let $a_0=p(p-1)+1$. Formula (DP.7), the low-weight
completeness just proved, and (DP.10) show that $W^0$
first occurs in weight $s+a_0$, where it is exactly
$RC$. The module $W^1$ first occurs in weight $s+a_0-1$,
where it is exactly $Rx$. At the first weight there is
no incoming relation in $P^0$, since the smallest
weight in $P^{-1}$ is $s+p^2>s+a_0$.
Consequently the entire arity-two degree-one lattice
in weight $2s+2a_0-1$ is
\[
                                 R\,Cx.             \tag{DP.15}
\]
Any primitive of (DP.13) in that arity and weight would
therefore be $\lambda Cx$ for some $\lambda\in R$.
Equation (DP.14) and $\xi\ne0$ require $p\lambda=1$,
which has no solution in $R$.
Arity is preserved by the differential, so a primitive
of another arity cannot alter this conclusion.
This proves exact order $p$ in $H^2(E)$.

We check the entire saturation. Form the algebraic
based closure of $E$ inside $E_{\mathbb Q}$ using its
augmentation $E\to S$: adjoin all rational zero-forms
with derivative already integral and integral augmentation,
close under PD products, and repeat countably many times.
Every stage is arity graded, since the derivative condition
separates by arity and the PD addition rule preserves
this grading. Arity zero remains exactly $S$.
No arity-one zero-form is added: for such an element
$f$, its derivative is a cycle in $W^1$; $H^1(W)=0$
supplies $b\in W^0$ with $db=df$, and
$H^0(W_{\mathbb Q})=0$ then gives $f=b$.
The same argument applies at every later stage.
All rational arity-two zero-forms have weight at least
$2s+2a_0$, one greater than (DP.15).
Products and PD operations cannot decrease weight or
arity. Hence every cochain group of arity two in the
witness weight is unchanged throughout the closure
and its union. The class (DP.13) retains exact order $p$.

We now give the actual twisted simplicial construction
announced in the theorem. Take (DP.1), tensor the parameter
of every $W$ generator with $K(2)$, and form the free
relative PD algebra over $X$. On the top generators use
\[
 du_0(t_0)=\varepsilon(t_0)z_0,\qquad
 dv_0(t_0)=p u_0(t_0)-\varepsilon(t_0)\eta_0,
 \tag{DP.16}
\]
where $\varepsilon:Q_0\to I_1$ is the displayed resolution
augmentation; on higher-resolution summands use (DP.9).
For the nonlinear twist, its value on the normalized
generator of $K(2)$ is the displayed polynomial, and
its other values are its linear Dold--Kan extension.
All faces of that value are zero, so this is an actual
simplicial linear map. The equality $d\eta_0=pz_0$ and
$\varepsilon\partial=0$ verify $d^2=0$ and every bottom
relation in (DP.10). Thus these formulas define an
actual strict simplicial PD DGA. Take its torsionfree
image in its rationalization, and call it $Y^{\rm tf}$.
There is an actual projection to $G_R$ killing the old
identity disk and all resolved generators; (DP.16)
shows that it respects the differential, and it factors
through the torsionfree image.

In the single old disk, the derivative
$d(x_0\gamma_{n-1}(y_0))=n\gamma_n(y_0)$ gives
$H^{2p}=R/p$ on $\gamma_p(y_0)$.
The explicit coefficient contraction of $(G_R)_2$
then gives the assertion about $z_0$ in $X_2$.
Rationally replace $u_0$ in (DP.16) by
$u_0-\varepsilon\eta_0/p$. This gives the untwisted
linear complex (DP.9). Its actual contraction descends
to (DP.10); extend it as a derivation on the rational
free algebra and divide its commutator with $d$ by
the positive arity. The commutator is multiplication
by that arity on generators and therefore on monomials.
This contracts every positive-arity summand.
It proves the rational coefficient comparison in each
simplex. Every finite monomial has finite arity, so
no completion or infinite-sum argument is involved.
The projection to $G_R$, together with that rational
comparison, identifies integral $H^0$ with $S$:
a closed degree-zero element is rationally a coefficient,
and its image in $G_R$ makes that coefficient integral.
The parameter $K(2)$ vanishes in simplices zero and one,
which proves the stated equalities there. The weights
in (DP.1) and (DP.9) also verify the support condition.
A-SAT01 therefore applies at cutoff $s+1$ with the
specified torsionfree target $G_R$.

In simplex two there is an actual augmented PD cochain map
\[
 Y_2^{\rm tf}\longrightarrow E:
 \quad (G_R)_2\longrightarrow S\text{ by vertex evaluation},
 \quad x_0,y_0\longmapsto0,
 \quad W\longmapsto W.                             \tag{DP.17}
\]
Its composite with $E\to S$ is exactly the source's
vertex-zero map. Consequently it extends at each based
saturation stage to the algebraic closure of $E$:
the image of a newly adjoined rational zero-form has
integral derivative at the preceding stage and has
integral augmentation. Products and divided powers
then extend this very map, and the filtered union
retains it. The polynomials (DP.13)--(DP.14) are still
identities in the twisted source, since their displayed
linear differentials are unchanged.
Their images under (DP.17) are the already detected
nonzero class. This proves the exact order claimed
in (DP.2) after the entire source saturation.

Finally the displayed class itself has an actual next
repair. Its value and $Cx$ on the normalized two-simplex
have all faces zero, so define actual $K(2)$ maps.
Adjoin a one-form $U$ and a zero-form $V$ on that
parameter with
\[
                      dU=\xi,\qquad dV=pU-Cx.
 \tag{DP.18}
\]
Equation (DP.14) verifies the differential identity.
The projection to $G_R$ extends by $U,V\mapsto0$.
The parameter still leaves simplices zero and one fixed,
and over $\mathbb Q$ the substitution $U'=U-Cx/p$
identifies the new pair with a PD identity disk
$V\mapsto pU'$. In particular the class has been killed
by an actual map-preserving repair. We make no assertion
here that its subsequent integral repairs converge to
an admissible theory.
\end{proof}

The independent certificate
\texttt{code/A\_cosimplicial\_deep\_tail\_certificate.py}
checks the complete coefficient lattices through $p^2$
for $p=3,5$, including the generator (DP.8), using
$54$ and $250$ exact local-lattice checks respectively.
Its recorded runtime is $0.09312325301289093$ seconds.
It verifies the displayed minimal weights, the entire
primitive lattice (DP.15), and the six coefficients
of (DP.14). The all-prime assertion is established by
the proof above, not extrapolated from those two checks.
